\documentclass[11pt]{article}
% Shared preamble for the printable mini lessons (1.4 onward).
%   \documentclass[11pt]{article}
%   \input{../preamble}
% Build the key with:  pdflatex -jobname=<name>_KEY "\def\iskey{}\input{<name>}"
% Figures shared with the web page live in figures/Unit_X/<name>.tex and are
% pulled in with \fig{Unit_X/<name>}; the standalone package drops their own
% preamble, so their macros come from figures/preamble.tex, loaded here.
% Shared preamble for web figures.  Each figure is a standalone .tex that
% \input's this file; scripts/make-figures.sh turns it into an SVG.
\usepackage[T1]{fontenc}
\usepackage{helvet}
\renewcommand{\familydefault}{\sfdefault}
\usepackage{tikz}
\usetikzlibrary{decorations.pathreplacing}
% axis numbers in the text font, not math italic
\pgfkeys{/pgf/number format/.cd,assume math mode=true,1000 sep={}}

% dot plot: \dotplot{min}{max}{value/count,...}{axis label}
\newcommand{\dotplot}[4]{%
\begin{tikzpicture}[x=0.9cm,y=0.33cm]
  \draw[thick] (#1-0.5,0) -- (#2+0.5,0);
  \foreach \t in {#1,...,#2} {\draw (\t,0.25) -- (\t,-0.25) node[below] {\t};}
  \foreach \v/\c in {#3} {\ifnum\c>0 \foreach \k in {1,...,\c} {\fill (\v,\k*1.15) circle (0.16cm);}\fi}
  \node at ({(#1+#2)/2},-2.3) {#4};
\end{tikzpicture}}

% shape sketch histogram (12 bins): \shapehist{heights}{label}
\newcommand{\shapehist}[2]{%
  \foreach \h [count=\i from 0] in {#1} \draw[fill=gray!30] (\i,0) rectangle (\i+1,\h);
  \draw (-0.3,0) -- (12.3,0);
  \node[font=\bfseries] at (6,-2.6) {#2};}
\newcommand{\meanmark}[1]{\fill (#1,0) -- ++(-0.3,-0.75) -- ++(0.6,0) -- cycle;}
\newcommand{\medmark}[2]{\draw[very thick,dashed] (#1,0) -- (#1,#2);}

% number line with minor ticks: \nline{min}{max}{label step}{axis label}
\newcommand{\nline}[4]{%
  \draw[thick] (#1-0.5,0) -- (#2+0.5,0);
  \foreach \t in {#1,...,#2} \draw (\t,0.12) -- (\t,-0.12);
  \pgfmathtruncatemacro{\nxt}{#1+#3}
  \foreach \t in {#1,\nxt,...,#2} \draw (\t,0.2) -- (\t,-0.2) node[below] {\t};
  \node at ({(#1+#2)/2},-0.95) {#4};}
% box plot above the line: \bplot{min}{Q1}{med}{Q3}{max}
\newcommand{\bplot}[5]{%
  \draw[thick] (#1,1.0) -- (#2,1.0);
  \draw[thick] (#4,1.0) -- (#5,1.0);
  \draw[thick] (#1,0.75) -- (#1,1.25);
  \draw[thick] (#5,0.75) -- (#5,1.25);
  \draw[thick,fill=gray!20] (#2,0.55) rectangle (#4,1.45);
  \draw[very thick] (#3,0.55) -- (#3,1.45);}

% ---- added for 1.4 onward ----
% axis on y=0: \axis{min}{max}{label step}{tick step}{axis label}
\newcommand{\axis}[5]{%
  \draw[thick] (#1,0) -- (#2,0);
  \pgfmathsetmacro{\axT}{#1+#4}\foreach \t in {#1,\axT,...,#2} \draw (\t,0.1) -- (\t,-0.1);
  \pgfmathsetmacro{\axL}{#1+#3}\foreach \t in {#1,\axL,...,#2}
    \draw (\t,0.18) -- (\t,-0.18) node[below] {\pgfmathprintnumber{\t}};
  \node at ({(#1+#2)/2},-0.95) {#5};}
% box plot centred at height y: \bplotat{y}{min}{Q1}{med}{Q3}{max}
\newcommand{\bplotat}[6]{%
  \draw[thick] (#2,#1) -- (#3,#1);
  \draw[thick] (#5,#1) -- (#6,#1);
  \draw[thick] (#2,#1-0.22) -- (#2,#1+0.22);
  \draw[thick] (#6,#1-0.22) -- (#6,#1+0.22);
  \draw[thick,fill=gray!20] (#3,#1-0.42) rectangle (#5,#1+0.42);
  \draw[very thick] (#4,#1-0.42) -- (#4,#1+0.42);}
% an outlier dot: \outlier{x}{y}
\newcommand{\outlier}[2]{\fill (#1,#2) circle (2.4pt);}

\usepackage[margin=0.75in]{geometry}
\usepackage{array}
\usepackage{enumitem}
\usepackage{amsmath}
\usepackage[helvet]{sfmath}  % math in the same sans font as the text
\usepackage{standalone}
\pagestyle{empty}
\setlength{\parindent}{0pt}
\setlist[enumerate,1]{leftmargin=1.5em,itemsep=1em}
\setlist[enumerate,2]{label=(\alph*),leftmargin=2em,itemsep=2.2em}

\newif\ifkey
\ifdefined\iskey \keytrue \fi
\colorlet{keyred}{red!75!black}
\newcommand{\ans}[1]{\ifkey\leavevmode{\color{keyred}#1}\fi}
% vertical work space on the student copy, less on the key
\newcommand{\work}[1]{\ifkey\vspace{0.6em}\else\vspace{#1}\fi}
\newcommand{\fig}[1]{\input{../../figures/#1}}
% Title line: \lessontitle{Quartiles, IQR, and Box Plots}
\newcommand{\lessontitle}[1]{{\Large\bfseries #1\ifkey\ \textcolor{keyred}{KEY}\fi}\hfill Name \underline{\hspace{2in}}\par\vspace{0.6em}}
\newcommand{\partner}{\par\vspace{0.4em}Work with a partner. Compare answers when you finish each problem.\par}

% Unit quiz: two questions per mini lesson, each headed by the lesson's
% objective.  The key keeps every box the same height as the student copy,
% so key page N lines up with quiz page N.
\newcommand{\lessonhead}[1]{\par\vspace{0.6em}{\large\bfseries #1}\par\vspace{0.25em}\hrule\vspace{0.6em}}
\newcommand{\goal}[1]{\textbf{#1}\par\vspace{0.35em}}
% answer box: same height on both copies, answer printed only on the key
\newcommand{\abox}[2]{\par\begin{minipage}[t][#1][t]{\linewidth}\raggedright\vspace{0.2em}\ans{#2}\end{minipage}\par}
% answer box inside a table cell: fixed height, so a key answer never makes a row taller
% (the \strut gives both copies the same first baseline)
\newcommand{\tbox}[2]{\begin{minipage}[t][#1][t]{\linewidth}\raggedright\strut\ans{#2}\end{minipage}}
% circled choices: \opt{A} prints plain; \optans{B} is circled in red on the key
\newcommand{\opt}[1]{\tikz[baseline=(c.base)]\node[rounded corners=0.6em,inner sep=2.5pt](c){#1};}
\newcommand{\optans}[1]{\ifkey\tikz[baseline=(c.base)]\node[rounded corners=0.6em,inner sep=2.5pt,draw=keyred,thick,text=keyred](c){#1};\else\opt{#1}\fi}
\setlist[enumerate,2]{label=(\alph*),leftmargin=2em,itemsep=0.3em}
\setlist[enumerate,1]{leftmargin=1.5em,itemsep=1em}
\raggedright\hyphenpenalty=10000\exhyphenpenalty=10000

\begin{document}

\lessontitle{Unit 5 Quiz: Study Design}

Show your work.

% ======================================================================
\lessonhead{5.1 - Samples, Studies, and Bias}
\begin{enumerate}

\item \goal{I name the population and the sample in a study, and I tell a sample survey, an observational study, and an experiment apart.}
A music app has 60,000 users. To learn how users feel about a new design, the app sends a survey to 500 randomly chosen users.

\begin{minipage}[t]{0.48\linewidth}\vspace{0pt}\raggedright
\begin{enumerate}
  \item What is the population?
  \abox{1.8em}{All 60,000 users of the app}
\end{enumerate}
\end{minipage}\hfill
\begin{minipage}[t]{0.48\linewidth}\vspace{0pt}\raggedright
\begin{enumerate}[start=2]
  \item What is the sample?
  \abox{1.8em}{The 500 users who got the survey}
\end{enumerate}
\end{minipage}

Is each study a sample survey, an observational study, or an experiment? Explain.

\vspace{0.2em}
\renewcommand{\arraystretch}{1.3}
\begin{tabular}{|>{\raggedright\arraybackslash}p{2.75in}|>{\raggedright\arraybackslash}p{1.3in}|>{\raggedright\arraybackslash}p{2.15in}|}\hline
\textbf{Study} & \textbf{Type of study} & \textbf{Explain}\\ \hline
(c) \; The music app study above & \tbox{2.6em}{Sample survey} & \tbox{2.6em}{It asks a sample questions. It changes nothing.}\\ \hline
(d) \; A researcher counts the birds that visit 20 yards with trees and 20 yards with no trees. & \tbox{3.4em}{Observational study} & \tbox{3.4em}{She only watches and counts. She does not change any trees.}\\ \hline
(e) \; A teacher randomly chooses half of her students to read a story on paper. The others read it on a screen. All take the same quiz. & \tbox{3.8em}{Experiment} & \tbox{3.8em}{She assigns a treatment (paper or screen) and measures the scores.}\\ \hline
\end{tabular}

\item \goal{I name the type of bias in a sample or a question, and I decide if results can be generalized to a population.}
Name the type of bias. Explain how it could change the results.
\begin{enumerate}
  \item A TV show asks viewers to go to its website and vote: should the show move to a later time?
  \abox{3em}{Voluntary response bias. Viewers choose to vote. People with strong feelings vote most, so the result may not match all viewers.}
  \item A survey asks, ``Should the city waste money on a new soccer field?''
  \abox{3em}{Response bias. The word ``waste'' pushes people to say no.}
  \item A gym emails a survey to 400 randomly chosen members. Only 90 members answer.
  \abox{3em}{Nonresponse bias. The 90 who answer may be different from the 310 who do not. For example, they may be the happiest or the most upset members.}
  \item Lina asks the 25 students in her computer club how many hours a day they use their phones. Can she generalize her results to all students at her school? Explain.
  \abox{3em}{No. This is a convenience sample, not a random sample. Computer club members may use screens more than other students, so the sample may not represent the school.}
\end{enumerate}

\end{enumerate}

% ======================================================================
\newpage
\lessonhead{5.2 - Sampling Methods}
\begin{enumerate}[resume]

\item \goal{I name and describe simple random, stratified, cluster, systematic, and convenience samples, and I tell stratified and cluster samples apart.}
A library has 1,500 members. It wants to survey 60 of them about new hours. Name each sampling method.

\vspace{0.2em}
\renewcommand{\arraystretch}{1.35}
\begin{tabular}{|>{\raggedright\arraybackslash}p{4.45in}|>{\raggedright\arraybackslash}p{1.75in}|}\hline
\textbf{How the library chooses 60 members} & \textbf{Method}\\ \hline
(a) \; A computer randomly picks 60 member numbers. & \ans{Simple random sample}\\ \hline
(b) \; A random member from the first 25, then every 25th member. & \ans{Systematic}\\ \hline
(c) \; 20 randomly chosen children, 20 teens, and 20 adults. & \ans{Stratified}\\ \hline
(d) \; The first 60 members who walk in on Saturday. & \ans{Convenience}\\ \hline
\end{tabular}

\vspace{0.4em}
\begin{enumerate}[start=5]
  \item A city randomly selects 4 of its 40 apartment buildings. It surveys every person who lives in those 4 buildings. Which sampling method does the city use? Circle the letter.

  \vspace{0.2em}
  \begin{tabular}{@{}l@{\hspace{0.4em}}l@{\hspace{3em}}l@{\hspace{0.4em}}l@{\hspace{3em}}l@{\hspace{0.4em}}l@{\hspace{3em}}l@{\hspace{0.4em}}l}
  \opt{A} & Simple random & \opt{B} & Stratified & \optans{C} & Cluster & \opt{D} & Systematic\\
  \end{tabular}
  \vspace{0.3em}
  \item Explain how a stratified sample of the people in the 40 buildings would be different from the city's sample.
  \abox{3.2em}{A stratified sample takes a few people, chosen at random, from \textbf{every} building. The city's cluster sample takes \textbf{everyone} from a \textbf{few} buildings.}
\end{enumerate}

\item \goal{I use a random digit table to choose a sample, and I choose a good sampling method for a situation.}
A coach has 40 players, numbered 01 to 40. Use this row of random digits to choose 4 players.
\begin{center}\large\texttt{\bfseries 34603 36904 52041 50165 61689 44176 49590 03997}\end{center}
\begin{enumerate}
  \item Write the pairs you read. Then write the 4 players you choose.
  \abox{3em}{Pairs: \textbf{34}, 60, \textbf{33}, 69, \textbf{04}, 52, 04, \textbf{15}. \; Skip 60, 69, 52 (over 40) and the second 04 (a repeat).\\ The sample is players 34, 33, 04, and 15.}
\end{enumerate}

A school has 900 students. 300 ride the bus, and 600 do not. The principal wants to ask 90 students about a new bus schedule.
\begin{enumerate}[start=2]
  \item Describe a good sampling method. Name it, and explain why it is a good choice.
  \abox{4.2em}{Stratified sample. Split the students into bus riders and others. Randomly pick 30 bus riders and 60 others (10\% of each). Bus riders may feel differently about the schedule, so both groups are in the sample in the right numbers. (Also accept an SRS of 90 with a reason.)}
  \item The principal asks the first 90 students who get off the buses. What is wrong with this method?
  \abox{3em}{It is a convenience sample. Only bus riders are asked. The 600 students who do not ride the bus are left out, so the results are biased.}
\end{enumerate}

\end{enumerate}

% ======================================================================
\newpage
\lessonhead{5.3 - Experiments}
\begin{enumerate}[resume]

\item \goal{I name the subjects, treatments, explanatory variable, and response variable in an experiment, and I explain why random assignment matters.}
A track coach wants to know if a new warm-up helps runners run faster. The 24 runners on the team are randomly split into two groups of 12. One group does the new warm-up. The other group does the usual warm-up. Then each runner runs 400 meters, and the coach records the time.

\vspace{0.4em}
\renewcommand{\arraystretch}{1.5}
\begin{tabular}{|>{\raggedright\arraybackslash}p{2.1in}|>{\raggedright\arraybackslash}p{4.1in}|}\hline
(a) \; Subjects & \ans{The 24 runners on the team}\\ \hline
(b) \; Treatments & \ans{The new warm-up and the usual warm-up}\\ \hline
(c) \; Explanatory variable & \ans{The type of warm-up}\\ \hline
(d) \; Response variable & \ans{Each runner's time for 400 meters}\\ \hline
\end{tabular}

\vspace{0.4em}
\begin{enumerate}[start=5]
  \item Why should the coach split the runners into groups at random? Explain.
  \abox{3.4em}{So the two groups are alike. The fastest runners do not all end up in one group. Then a difference in times comes from the warm-up, not from the runners.}
  \item Instead, the coach lets each runner choose a warm-up. Why is this a problem?
  \abox{3.4em}{The groups may not be alike. The fastest runners may all choose the new warm-up. Then we cannot tell if the warm-up or the runners caused the faster times.}
\end{enumerate}

\item \goal{I explain control groups, placebos, and blinding, and I describe a completely randomized design and a matched pairs design.}
A company says its new ``focus gum'' helps students remember. 40 student volunteers are randomly split into two groups of 20. One group chews focus gum. The other chews plain gum that looks and tastes the same. Then all students take the same memory quiz, and a teacher grades it.
\begin{enumerate}
  \item Which group is the control group?
  \abox{1.6em}{The 20 students who chew the plain gum}
  \item The plain gum is a placebo. Why give the control group a placebo, and not nothing at all?
  \abox{3.4em}{Students may do better just from \textit{thinking} they got focus gum (the placebo effect). With a placebo, both groups think so. Then an extra difference comes from the gum itself.}
  \item The teacher who grades the quizzes does not know which gum each student chewed. Why?
  \abox{3em}{If she knew, she might, without meaning to, grade the focus gum group more kindly. Not knowing keeps the grading fair.}
  \item This experiment is a completely randomized design. Describe how to test the gum with a matched pairs design instead.
  \abox{3.4em}{Each student takes two memory quizzes of the same difficulty: once with focus gum and once with plain gum. A coin flip decides which gum comes first. Compare each student's two scores.}
\end{enumerate}

\end{enumerate}

% ======================================================================
\newpage
\lessonhead{5.4 - What a Study Can Prove}
\begin{enumerate}[resume]

\item \goal{I decide if a study's results can be generalized to a population, and I decide if a study can show cause and effect.}
For each study, can we conclude cause? Can we generalize to the population? Explain.

\vspace{0.2em}
\renewcommand{\arraystretch}{1.3}
\begin{tabular}{|>{\raggedright\arraybackslash}p{2.85in}|>{\raggedright\arraybackslash}p{1.65in}|>{\raggedright\arraybackslash}p{1.7in}|}\hline
\textbf{Study} & \textbf{Cause? Why?} & \textbf{Generalize? Why?}\\ \hline
(a) \; Researchers randomly choose 600 teens in Missouri. Teens who spend more time outside sleep more. & \tbox{3.6em}{No. Observational study, no random assignment.} & \tbox{3.6em}{Yes, to Missouri teens. The sample is random.}\\ \hline
(b) \; 60 volunteers from a typing class are randomly assigned to Game A or Game B. The Game A group types faster. & \tbox{3.6em}{Yes, for these 60 (random assignment).} & \tbox{3.6em}{No. Volunteers are not a random sample.}\\ \hline
\end{tabular}

\vspace{0.3em}
\begin{enumerate}[start=3]
  \item A researcher randomly selects 200 students from a school district. She randomly assigns half to study with flash cards and half with a video. The flash card group scores higher, and the difference is statistically significant. Which conclusion is best supported?

  \vspace{0.2em}
  \begin{tabular}{@{}l>{\raggedright\arraybackslash}p{5.6in}}
  \opt{A} & Studying with flash cards causes higher scores for all students in the United States.\\
  \optans{B} & Studying with flash cards causes higher scores for students in this school district.\\
  \opt{C} & Flash cards are associated with higher scores, but the study cannot show cause.\\
  \opt{D} & No conclusion can be made, because not every student in the district was studied.\\
  \end{tabular}
\end{enumerate}

\item \goal{I use a simulation to decide if a result is statistically significant, and I find mistakes in conclusions about studies.}
A teacher tests a new vocabulary game. 30 students are randomly split into two groups of 15. The game group's mean is 4 points higher. Each dot is the difference between the group means in one random shuffle of the 30 scores.

\vspace{0.3em}
\begin{minipage}[t]{0.47\linewidth}\vspace{0pt}\raggedright
\begin{tikzpicture}[x=0.55cm,y=1cm]
  \axis{-6}{6}{1}{1}{Difference between the group means}
  \foreach \v/\c in {-6/1,-5/2,-4/5,-3/9,-2/12,-1/15,0/16,1/16,2/14,3/7,4/2,5/1}
    {\foreach \k in {1,...,\c} {\fill (\v,\k*0.18) circle (0.07cm);}}
  \node[font=\small,anchor=west] at (2.4,2.6) {100 shuffles};
\end{tikzpicture}
\end{minipage}\hfill
\begin{minipage}[t]{0.5\linewidth}\vspace{0pt}\raggedright
\begin{enumerate}
  \item How many shuffles gave a difference of 4 or more?
  \abox{1.6em}{$2 + 1 = 3$ shuffles}
  \item Is a 4-point difference statistically significant? Explain.
  \abox{3.8em}{Yes. Chance gave a difference this big only 3 of 100 times (3\%). That is less than 5\%.}
\end{enumerate}
\end{minipage}

\vspace{0.3em}
\begin{minipage}[t]{0.48\linewidth}\vspace{0pt}\raggedright
\begin{enumerate}[start=3]
  \item Find the mistake. In another class, the game group scored 2 points higher. The teacher says, ``The game works.''
  \abox{4.4em}{Chance gave 2 or more $14 + 7 + 2 + 1 = 24$ times out of 100 (24\%). That is common. 2 points is not statistically significant.}
\end{enumerate}
\end{minipage}\hfill
\begin{minipage}[t]{0.48\linewidth}\vspace{0pt}\raggedright
\begin{enumerate}[start=4]
  \item Find the mistake. A survey of 500 randomly chosen adults finds that dog owners walk more. A headline says, ``Owning a dog makes people walk more.''
  \abox{4.4em}{A survey shows an association, not a cause. People who already like to walk may choose to own a dog.}
\end{enumerate}
\end{minipage}

\end{enumerate}

\end{document}
